\documentclass[10pt,a4paper]{amsart}
\usepackage[margin=19mm]{geometry}
\usepackage{amsmath,amssymb,mathtools}
\usepackage[scr=rsfs,cal=euler]{mathalfa}
\usepackage{xfrac}
\usepackage[unicode,hidelinks,bookmarks=false]{hyperref}
\newcommand{\ie}{i.e.\ }
\newcommand{\cf}{cf.\ }
\newcommand{\resp}{respectively\ }
\newcommand{\Subsection}{\subsection}
\providecommand{\nobreakdash}{\nobreak\hbox{-}}
\setlength{\emergencystretch}{2em}
\multlinegap0pt
\usepackage{bm}
\usepackage{bbm}
\usepackage{dsfont}
\usepackage{xspace}
\usepackage{cancel}
\usepackage{mathtools}
\usepackage{amsthm}
\makeatletter
\@ifundefined{theorem}{\newtheorem{theorem}{Theorem}}{}
\@ifundefined{proposition}{\newtheorem{proposition}[theorem]{Proposition}}{}
\makeatother
\DeclareMathOperator{\tr}{tr}
\newcommand{\cW}{\mathcal{W}}
\newcommand{\dd}{\mathrm{d}}
\title[Topological Recursion and Quantum Path Signatures]{Topological Recursion and Quantum Path Signatures}
\author{In\^es Aniceto, Thomas Cass, Samuel Crew}
\date{Source: arXiv:2608.02861v1 (CC BY 4.0)}
\begin{document}
\maketitle

\noindent\emph{Source and licence.} Topological Recursion and Quantum Path Signatures. Version arXiv:2608.02861v1 (CC BY 4.0), distributed under the Creative Commons Attribution 4.0 International licence, as recorded in the arXiv licence metadata for this exact version and shown on the abstract page https://arxiv.org/abs/2608.02861v1. Full rights notice: licenses/arxiv-2608.02861-CC-BY-4.0.txt.\par\smallskip

\noindent\textit{Editorial scope.} Counted unit: the Proposition whose source label is \texttt{prop:W2-loop} in the article (source0.tex lines 867--878, source label \texttt{prop:W2-loop}), delivered together with the article's own complete original proof of it (source0.tex lines 880--936, one \texttt{proof} environment of 57 lines). No mathematics is written, repaired or re-derived: statement and proof are the article's own text line for line. Required context retained uncounted: eq:two-trace (lines 853--855). Every definition, display and label of those lines is retained unchanged. Omitted from this package and not redistributed: 1--852, 856--866, 879--879, 937--960. Nothing inside a retained excerpt is omitted or altered: every retained body is delivered exactly as the article's lines are written, so no omission receipt of any kind accompanies this package. Citations: the retained excerpts carry no citation command at all, so no bibliography entry of the article is transcribed and no reference is redistributed. Reference adaptation: none was needed and none was made; every reference inside the delivered unit resolves inside it, so no unresolved reference or citation is produced. Printed slips kept literal: no printed word, symbol, hypothesis or number of the counted statement or of its proof was altered. Layout and class: the article's own class is \texttt{article} with packages including amsmath, amssymb, amsthm, geometry, hyperref, mathtools, enumitem, graphicx, authblk; the wrapper is the lane's pinned \texttt{amsart} editorial wrapper at 10pt on A4, which restates the theorem-like environments the retained text needs and adds the article's own macro definitions that the retained text needs (\texttt{\textbackslash cW}, \texttt{\textbackslash dd}, \texttt{\textbackslash tr}), with extra wrapper packages (bm, bbm, dsfont, xspace, cancel, mathtools). The delivered numbering is the unit's own. Assets: no retained span contains a figure environment, a graphics-inclusion command, a PDF-page inclusion, a table or a TikZ picture; no image file is copied into this package, the package declares no asset dependency and no image file is redistributed with it. Licence: version arXiv:2608.02861v1 of this article is distributed under the Creative Commons Attribution 4.0 International licence, as recorded in the arXiv licence metadata for this exact version and shown on the abstract page https://arxiv.org/abs/2608.02861v1. A full rights notice is delivered at licenses/arxiv-2608.02861-CC-BY-4.0.txt. Characters: every retained excerpt is pure ASCII, so the delivered engine, XeLaTeX with its default Latin Modern text font, reports no missing character.
\begin{equation}\label{eq:two-trace}
  \bigl\langle\tr(A_I A_\mu)\,\tr A_J\bigr\rangle = \frac1N\sum_{I=K\mu L}\bigl\langle\tr A_K\,\tr A_L\,\tr A_J\bigr\rangle + \frac1N\sum_{J=P\mu Q}\bigl\langle\tr(A_I\,A_Q A_P)\bigr\rangle.
\end{equation}

\begin{proposition}[Loop equation for the two-point function]
\label{prop:W2-loop}
The planar two-point function $\cW_{0,2}(a,b;c,d)$ satisfies
\begin{equation}\label{eq:W2-loop}
  \cW_{0,2}=L_0\cW_{0,2}+\beta_0,
\end{equation}
where
\begin{equation}\label{eq:beta0}
\beta_0(a,b;c,d) = -\int_a^b\!\!\int_c^d \bigl\langle\dd\gamma(u),\dd\gamma(v)\bigr\rangle \cW_0\bigl( \gamma|_{[a,u]}* \gamma|_{[v,d]}* \gamma|_{[c,v]} \bigr).
\end{equation}
Hence $\cW_{0,2}$ is determined by $\cW_0$.
\end{proposition}

\begin{proof}
The connected correlator is
\begin{equation}
\langle X;Y\rangle_{\mathrm c}
=
\langle XY\rangle-\langle X\rangle\langle Y\rangle
\end{equation}
Subtracting the product of the one-trace identity with $\langle\tr A_J\rangle$ from \eqref{eq:two-trace} gives
\begin{equation}\label{eq:two-trace-conn}
\begin{aligned}
\bigl\langle\tr(A_I A_\mu);\tr A_J\bigr\rangle_{\mathrm c} =& \frac{1}{N}\sum_{I=K\mu L} \bigl\langle\tr A_K\,\tr A_L;\tr A_J\bigr\rangle_{\mathrm c}
\\ &+ \frac{1}{N}\sum_{J=P\mu Q} \bigl\langle\tr(A_I A_Q A_P)\bigr\rangle.
\end{aligned}
\end{equation}
At leading order in $N$,
\begin{equation}\label{eq:three-trace-factorisation}
\begin{aligned}
\bigl\langle\tr A_K\,\tr A_L;\tr A_J\bigr\rangle_{\mathrm c}
=& \langle\tr A_K\rangle \bigl\langle\tr A_L;\tr A_J\bigr\rangle_{\mathrm c}
\\ &+ \bigl\langle\tr A_K;\tr A_J\bigr\rangle_{\mathrm c} \langle\tr A_L\rangle +O(N^{-1}).
\end{aligned}
\end{equation}
Define the planar connected moments by
\begin{equation}\label{eq:planar-connected-moments}
\begin{aligned}
M^{(0),\mathrm c}_{I;J}
&=
\lim_{N\to\infty}
\bigl\langle
\tr\bigl((iA)_I\bigr);
\tr\bigl((iA)_J\bigr)
\bigr\rangle_{\mathrm c}
\end{aligned}
\end{equation}
Using \eqref{eq:three-trace-factorisation} in \eqref{eq:two-trace-conn} gives
\begin{equation}\label{eq:two-trace-planar-recursion}
M^{(0),\mathrm c}_{I\mu;J} = -\sum_{I=K\mu L} \left( M_K M^{(0),\mathrm c}_{L;J} + M^{(0),\mathrm c}_{K;J}M_L \right) -\sum_{J=P\mu Q}M_{IQP}.
\end{equation}
The signs come from the factor $i^2=-1$ associated with each
contraction.

Expanding the two developments against their signatures gives
\begin{equation}\label{eq:W02-signature-expansion}
\cW_{0,2}(a,b;c,d) = \sum_{I,J} M^{(0),\mathrm c}_{I;J} S^I_{a,b}S^J_{c,d}.
\end{equation}
Multiply \eqref{eq:two-trace-planar-recursion} by $S^{I\mu}_{a,b}S^J_{c,d}$ and sum over $I,J,\mu$. The first sum gives
\begin{equation}
-\int_{a<u<r<b} \bigl\langle\dd\gamma(u),\dd\gamma(r)\bigr\rangle \left[ \cW_0(a,u)\cW_{0,2}(u,r;c,d) + \cW_{0,2}(a,u;c,d)\cW_0(u,r) \right],
\end{equation}
which is $L_0\cW_{0,2}$.

In the second sum, the occurrence of $\mu$ in $J=P\mu Q$ marks a point $v$ on the second boundary. Summing over $\mu$ and using Chen's identity gives
\begin{equation}
-\int_a^b\!\!\int_c^d \bigl\langle\dd\gamma(u),\dd\gamma(v)\bigr\rangle \cW_0\bigl( \gamma|_{[a,u]}* \gamma|_{[v,d]}*\gamma|_{[c,v]} \bigr),
\end{equation}
which is $\beta_0$. 
\end{proof}

\end{document}
